% GUIA DO PROFESSOR — Volume 5 (gerado por expandir_volume.py)

\chapter{Organização do Tempo de Ensino}

No 5º ano, a rotina diária (10–15 min de fônico, 10–15 min de Kumon, 30 min de leitura e escrita) estrutura o tempo. O quadro de horários distribui as unidades da sequência didática e as folhas Kumon por semana, respeitando o ritmo da turma.


\chapter{Materiais Necessários}

Materiais: livro do aluno, folhas Kumon, espelho pequeno, cartazes da Boquinha, fichas de palavras, caderno de pauta dupla, cronômetro e o Registro de Progresso. Cada aluno mantém uma pasta de folhas para revisão quinzenal.


\chapter{Estratégias de Diferenciação}

Usar as rotas 1A–1D: alunos que consolidaram avançam; alunos com dificuldade repetem folhas e recebem apoio da Boquinha e do espelho. A diferenciação é por evidência (Registro de Progresso), nunca por adivinhação.


\chapter{Registro de Progresso}

O Registro é o coração do método: cada folha corrigida entra no registro com data, acertos e erros; os checkpoints alimentam o quadro S/F/R/N. Revisão quinzenal orienta as rotas.


\chapter{Atividades Complementares}

Jogos com fichas de palavras, ditados interativos, caça-palavras e produção coletiva de textos. Atividades complementares reforçam o som-alvo sem virar lição de casa mecânica.


\chapter{Orientações para Famílias}

Pauta da reunião de pais: o que é o método, como ajudar em casa (5–10 min/dia de leitura compartilhada e ditado de palavras das unidades), como funciona o rastreio escolar (não é diagnóstico) e quando buscar avaliação profissional.


\chapter{Alinhamento com as Diretrizes Curriculares Nacionais (DCNC)}

As DCNC e a BNCC orientam a progressão do 5º ano: leitura e produção de gêneros do cotidiano, análise linguística e consciência fonológica avançada. O livro converge com o que a rede já ensina, sem sobreposição.


\chapter{Rotas Individualizadas e Neuroinclusão na Prática}

Estudos de caso: aluno com dislexia na família (rota 1A + repetição de planchetas), aluno com TDAH (pausas curtas, apoio visual, checklist de passos), aluno com altas habilidades (rota 1D, produção ampliada).


\chapter{O Papel do Professor no Rastreio Escolar}

O professor observa, registra e encaminha. NUNCA diagnostica. Os checkpoints de rastreio (SPEC-935-R211) e o guia de triagem ajudam a organizar a observação; o Volume Profissional de Sondagem e Rastreio é instrumento do profissional habilitado.


\chapter{Avaliação da Própria Prática}

Autoavaliação docente mensal: os alunos avançaram nas rotas? Quais folhas geraram erros repetidos? As reuniões pedagógicas usam o Registro de Progresso como evidência para ajustar o planejamento.


\chapter{Planejamento por Bimestre}

O 5º ano se organiza em quatro bimestres. O primeiro bimestre consolida a base do ano anterior e apresenta os dígrafos CH, LH e NH; o segundo trabalha QU/GU e os encontros consonantais; o terceiro foca as regras ortográficas RR, SS, S/Z e Ç; o quarto encerra com AM/AN, X, revisões e leitura expressiva de gêneros. Cada bimestre encerra com uma avaliação e um checkpoint de rastreio.


\chapter{Exemplos de Aula Passo a Passo}

Aula de 45 minutos do dígrafo CH: (1) aquecimento com a Boquinha e o espelho (5 min); (2) formação guiada das sílabas CHA–CHU com o quadro e as fichas (10 min); (3) leitura das palavras da unidade com apoio (10 min); (4) folha de prática 1, itens de reconhecimento (10 min); (5) fechamento com a autoavaliação oral (5 min). A lição de casa: reler o texto curto da unidade com um adulto.


\chapter{Projetos de Leitura e Biblioteca}

O projeto 'Minha Primeira Coleção' leva cada aluno a escolher 5 livros no ano, registrar o título e o autor no seu caderno de leitor e apresentar um resumo oral por trimestre. A biblioteca vira ambiente de fluência: leitura em voz alta em duplas, gravação de recitação e sarau de poesia no fim do ano. O rastreio observa, sem diagnosticar, quem ainda precisa de apoio para ler com ritmo e expressividade.


\chapter{Avaliando a Fluência na Rotina}

A fluência se mede pela precisão, ritmo e expressividade, não pela velocidade. Na rotina, use leituras de 1 minuto em textos conhecidos e desconhecidos; registre as palavras lidas corretamente; compare com a linha de base do início do ano. A meta do 2º ano é ler textos curtos com ritmo e entonação, com compreensão.


\chapter{Jogos com Sílabas e Palavras}

Jogos de baixo custo: dominó de sílabas (CHA-CHE-CHI), bingo de palavras-chave, corrida de leitura em duplas, caça-palavras das grafias e baralho de encontros consonantais. Cada jogo termina com registro escrito no caderno: a brincadeira sempre converte em escrita.


\chapter{Adaptações para Alunos com Suspeita de Dislexia, TDAH e TDL}

Para suspeita de dislexia: mais repetição nas folhas, leitura em voz alta com apoio, fontes maiores e tempo extra. Para TDAH: instruções curtas, checklist de passos e pausas programadas.


Para transtorno de linguagem: apoio da Boquinha diário, frases curtas no ditado e repetição com espelho. São adaptações pedagógicas; o diagnóstico é sempre do profissional habilitado.


\chapter{Integração com o Volume Profissional e o Caderno Motor}

O Volume Profissional traz a triagem inicial e as fichas de observação que identificam o aluno em risco; este volume aplica as rotas 1A–1D e as folhas Kumon; o Caderno Motor desenvolve a motricidade fina que antecede a letra cursiva. A sequência recomendada no ano: triagem no início, checkpoints a cada três unidades e reavaliação no fim de cada bimestre, sempre com a integralidade do programa, nunca de forma isolada.


\chapter{Rotina para Crianças com Atraso na Fala}

Quando a fala está em atraso, a consciência fonológica trava: a criança não diferencia sons que ela mesma não articula. A rotina recomendada: trabalho diário com a Boquinha no espelho, associação gesto-som de cada grafia nova, repetição das sílabas da unidade em voz alta e leitura em eco (o professor lê, a criança repete). A família deve ser orientada a estimular a fala em casa, e o professor deve sinalizar à coordenação quando a fala estiver comprometendo a aprendizagem, para encaminhamento pelo profissional habilitado.

